The verification run completed successfully on 12 September 2026 using
\texttt{leanprover/lean4:v4.30.0-rc2}. Run from the project directory:
\begin{Verbatim}[breaklines=true]
LEAN_NUM_THREADS=4 python3 scripts/verify.py
\end{Verbatim}
The script reported \texttt{verified}. It checked 1,485 upstream source
entries against their manifest hashes, checked nine dependency revisions
against the upstream lock file, recorded local source SHA-256 hashes,
and completed the build of both libraries (5,105 build jobs, including
replayed jobs). It then audited 255 listed theorem endpoints. The public
target-density theorem and counterexample dichotomy use only
\texttt{propext}, \texttt{Classical.choice}, and \texttt{Quot.sound}.
No audited endpoint used an axiom outside this standard set. The source
escape scan also passed; in particular the scanned proof sources contain
no executable \texttt{sorry} or additional axiom declaration.

The machine-readable record is \texttt{verification/local-report.json};
build and axiom logs are in the same directory. Every one of the 13 local
modules reproduced below was part of this build. The appendix contains
1,283 source lines and was compared byte-for-byte with those modules.
It relies on the project's pinned imported libraries, so compilation of
isolated snippets without those imports is not the verification procedure.

To regenerate this PDF from its output directory, run XeLaTeX twice:
\begin{Verbatim}[breaklines=true]
xelatex -interaction=nonstopmode -halt-on-error collatz_target_density_report.tex
xelatex -interaction=nonstopmode -halt-on-error collatz_target_density_report.tex
\end{Verbatim}
The adjacent \texttt{fonts/JuliaMono-Regular.ttf} supplies the Unicode
characters in the Lean listings.
